%=====================================================================
% FLOW: ABSTRACT (~200-230 words). Structured but unlabeled.
% Beats: (1) problem+stakes; (2) gap/question RQ2; (3) data+method;
%        (4) headline result (frontier + RQ2 sign + economic USD);
%        (5) contribution/benchmark. Write near-last; keep every number sourced.
%=====================================================================
\begin{abstract}
\noindent
\wnote{one paragraph, no citations, no undefined acronyms}
Liquidation cascades are the systemic-risk analogue of decentralized (DeFi)
lending: a price move forces collateral sales that push prices further, forcing
more liquidations. \todoitem{1 sentence on scale/recurrence — five episodes,
\$$\sim$X liquidated.} We ask whether such cascades can be \emph{forecast}
ahead of time, and in particular whether modeling the \emph{contagion network}
between positions, collateral assets, and venues carries predictive information
beyond aggregating independent position risks (RQ2).
We assemble \todoitem{the largest labeled cascade dataset for DeFi lending}
(Aave, Compound, Maker, and modular markets; Ethereum 2021--Feb 2026; five
systemic episodes, two with mempool coverage), reconstruct position health
factors, and adopt a \emph{pre-registered} cascade definition reconciled to
protocol-published statistics. We couple multiplex-network Hawkes processes with
temporal graph neural networks to produce calibrated $\Prob(\text{cascade within
}\horizon\text{ blocks})$ and severity, and evaluate under strict expanding-window
walk-forward with episode hold-outs.
\todoitem{Headline: network vs.\ aggregate AUPRC gain (or null) with block-bootstrap
CI; lead-time--precision frontier; \$$X$M protectable at $\langle h\rangle$-block
lead time.} We release the dataset, labels, and code as an open benchmark.
\end{abstract}
